% !TEX root = ../main.tex

\chapter{双喷嘴横向射流火焰实验历程}

本章介绍为获取横向射流火焰熄火边界而开展的实验探索过程，包括前置工况筛选及声激励传递特性的初步评估。

\section{前置实验与工况探索}

该部分是对正式实验的工况探索实验，目的是先对不同工况进行实验，观察不同工况下火焰的燃烧情况，挑选合适的工况进行下一步实验。截至目前，已完成多次不同工况下火焰燃烧情况的实验，对不同工况下的火焰燃烧情况进行记录。

在进行对二级横向射流火焰的研究之前，需要对一级燃烧室的多重扩散火焰的工况进行确定，由于燃烧器的性质，要使一级燃烧室产生的热流经过收缩段之后，能够点燃二级燃烧室通入的气体，因此需要对一级燃烧室的火焰燃烧情况以及所产生的热流温度进行探测，对一级燃烧室的火焰进行图像采集，并用热电偶记录其温度，如图~\ref{fig:34}所示为多个工况火焰燃烧的情况。本实验采用的气体为压缩空气和甲烷，当甲烷充分燃烧时，产生的火焰是淡蓝色的，图片中有黄色的火焰出现说明有燃烧不充分的现象存在，正式实验需要避免这些工况。

\begin{figure}[!htbp]
	\centering
	\includegraphics[width=0.7\textwidth]{image34.png}
	\caption{不同工况下多重扩散火焰燃烧情况}
	\label{fig:34}
\end{figure}

针对适合用于二级横向射流火焰燃烧的一级燃烧工况进行筛选，要求是能够产生纯淡蓝色的火焰，并且能产生1200~K--1500~K左右的温度，使用这些工况去对二级燃烧室横向射流火焰进行研究。目前主要使用的工况如表~\ref{tab:4}所示：

\begin{table}[!htbp]
	\centering
	\caption{现阶段实验工况表}
	\label{tab:4}
	\begin{tabular}{c|c}
		\hline
		甲烷流量（SLM） & 空气流量（SLM） \\
		\hline
		44 & 360 \\
		42 & 510 \\
		44 & 650 \\
		47 & 820 \\
		51 & 970 \\
		\hline
	\end{tabular}
\end{table}

由于在环境光较亮的情况下，淡蓝色火焰不易拍摄，二级燃烧室的气体先采用乙烯和空气进行现象观测，后续进行正式实验使用甲烷，图~\ref{fig:35}为使用乙烯，采用单喷嘴时所拍摄的二级燃烧室横向射流火焰的图片。

\begin{figure}[!htbp]
	\centering
	\includegraphics[width=0.7\textwidth]{image35.jpg}
	\caption{横向射流火焰燃烧情况（乙烯/单喷嘴）}
	\label{fig:35}
\end{figure}

6~mm双喷嘴加工好后，采用双喷嘴进行二级燃烧室的实验，实验工况如表~\ref{tab:5}所示，所得实验结果如图~\ref{fig:36}、~\ref{fig:37}所示。本次实验所使用的工况产生了黄色火焰居多的燃烧现象，这是燃烧不充分的表现，因此还需对二级燃烧室的实验工况进行调整，并且采用孔径更小的双喷嘴进行实验。

\begin{table}[!htbp]
	\centering
	\caption{二级燃烧室横向射流火焰工况表（乙烯/空气）}
	\label{tab:5}
	\begin{tabular}{c|c|c}
		\hline
		编号 & 乙烯流量（L/min） & 空气流量（L/min） \\
		\hline
		1 & 4  & 0  \\
		2 & 4  & 25 \\
		3 & 4  & 30 \\
		4 & 6  & 35 \\
		5 & 6  & 40 \\
		6 & 6  & 50 \\
		7 & 8  & 40 \\
		8 & 8  & 50 \\
		\hline
	\end{tabular}
\end{table}

\begin{figure}[!htbp]
	\centering
	\includegraphics[width=0.9\textwidth]{image36.jpg}
	\caption{横向射流火焰实验燃烧图1（6~mm双喷嘴）}
	\label{fig:36}
\end{figure}

\begin{figure}[!htbp]
	\centering
	\includegraphics[width=0.9\textwidth]{image21.jpg}
	\caption{横向射流火焰实验燃烧图2（6~mm双喷嘴）}
	\label{fig:37}
\end{figure}

该部分对一级燃烧室内的燃烧工况进行了探索，进行前置实验，并对使用工况进行筛选，并且已经完成对二级燃烧室内横向射流火焰的燃烧情况进行观察及拍照，为后续实验奠定坚实的基础，后续需改变二级燃烧室的双喷嘴的内径以及工况进行修改，进行后续实验。

\section{激励传递特性初步评估}

由于本项目需要使用集成的声激励脉动系统对气体进行激励，因此需要对其激励效果进行评估，于是借用学院的热膜风速仪进行测试在激励情况下速度波动的幅值。美国TSI IFA300热膜风速仪的主要技术参数如表~\ref{tab:热膜风速仪}所示。

\begin{table}[!htbp]
	\centering
	\caption{美国TSI IFA300热膜风速仪主要技术参数}
	\label{tab:热膜风速仪}
	\begin{tabular}{c|c}
		\hline
		参数名称 & 技术指标 \\
		\hline
		放大器输入漂移 & 0.3 μV/℃ \\
		等效放大器输入噪声 & 1.7 mV/Hz \\
		频率响应 & 300 kHz（无需方波试验） \\
		精度 & 0.1\% 或 0.01 Ω \\
		工作电阻 & 2--80 Ω \\
		最大探针电流 & 1.6 A \\
		电缆长度 & 5 m 或 30 m \\
		输出 & 桥电压到 11.5 V \\
		尺寸 & 17.8 cm × 48 cm × 41 cm \\
		通道数 & 每箱8通道，共16通道 \\
		电源 & 100/120/220/240 VAC，50--60 Hz，不需转换开关 \\
		\hline
		\multicolumn{2}{c}{信号适配器} \\
		\hline
		偏置 & 0--10 V，10 mV，0.15\% 精度 \\
		增益 & 1--1000（最大1M带宽），0.15\% 精度 \\
		低通滤波 & 12 Hz 到 1 MHz，13档，-60 dB/倍频程 \\
		高通滤波 & 0.1，1，10 Hz，-60 dB/倍频程 \\
		输出（模拟） & -5 到 +5 V \\
		最高采样率 & 1 MHz \\
		A/D板 & 4/8通道，12 bit，-5～+5 V，每通道最高采样率：714.3 kHz / 1 MHz \\
		\hline
	\end{tabular}
\end{table}

将气管连接到声激励喷嘴之后，再从声激励喷嘴的出口处连接到热膜风速仪进行测试，通过设置流量以及调节信号发生器所设置的频率和幅值，记录示波器所显示的波形及改变输入信号时的变化。

测试结果如图~\ref{fig:39}所示。

\begin{figure}[!htbp]
	\centering
	\includegraphics[width=0.8\textwidth]{image38.jpg}
	\caption{示波器采集界面}
	\label{fig:39}
\end{figure}

实验现象是，示波器无法正确显示出所设置的波形，不论如何调整信号发生器所输入的信号，示波器都无法显示出波形，只有在调整输入气体流量时，示波器上的波形会随着流量的变大变小而上下跳动。热膜风速仪无法响应高频信号，未能测得有效数据。

需要对测试进行改进，计划采用双色PIV速度场数据直接计算脉动幅值，或借用高频PIV/热线风速仪重新测试。